\section{Analysis}
\label{sec:analysis}

\begin{figure}[t]
  \centering
  \figorpending{qualitative_transport.pdf}{5.0cm}{Qualitative grid (DROID and surgical test episodes):
  observed frame, future ground truth, learned transport field (arrows = expected source displacement per
  patch), gate map $g_k$, optical-flow reference (RAFT), and per-patch error maps of AR vs.\ \ours{},
  decoded to RGB for steps $k\in\{1,5,10\}$.}
  \caption{\textbf{The transport field is a readable forecast of motion.} For each held-out example
  (rows), the arrows show where each predicted patch draws its features from, and colour shows the gate.
  Transport concentrates on the moving arm, gripper and instruments. The static background keeps the
  identity path, and newly revealed regions fall back to the correction term. The rightmost columns
  compare decoded forecasts and error maps with the strongest baseline.}
  \label{fig:qualitative}
\end{figure}

\paragraph{What does the transport field learn? (Q4)} We convert $\pi_{k,i}$ into an expected source
displacement per patch and compare it, at patch resolution, with optical flow estimated by RAFT
\citep{teed2020raft} between the observed and future frames (\cref{fig:qualitative}). \pend{}

\paragraph{Where do the gains come from?} We bin test patches by the magnitude of their true feature
change and compare the per-bin error of \ours{} and Direct (\cref{app:analysis}). \pend{}

\begin{figure}[t]
  \centering
  \begin{minipage}[t]{0.49\linewidth}\centering
    \figorpending{tradeoff.pdf}{3.9cm}{Forecast error (x) vs.\ action ranking accuracy (y) for every
    predictor and ablation; marker size = planning success.}
  \end{minipage}\hfill
  \begin{minipage}[t]{0.49\linewidth}\centering
    \figorpending{gain_vs_motion.pdf}{3.9cm}{Relative MSE gain of \ours{} over Direct, binned by true
    per-patch feature change, DROID and Hamlyn test.}
  \end{minipage}
  \caption{\textbf{Left: accuracy vs.\ action sensitivity.} Anchored predictors can reach low open-loop
  error while their forecasts react weakly to the actions. The action-contrastive objective moves \ours{}
  toward the upper-left. \textbf{Right: gains concentrate where the scene moves.}}
  \label{fig:tradeoff}
\end{figure}

\paragraph{Accuracy versus action sensitivity.} A forecast that stays close to $\Z_0$ can score
well on average while barely responding to the commanded actions, and a planner then cannot rank
candidate action sequences. We measure this directly as the rate at which true future actions yield
lower error than actions from another episode (\cref{fig:tradeoff}, left). \pend{}

\begin{table}[t]
\centering
\caption{\textbf{Efficiency} for a 10-step forecast of a $16{\times}16$ grid on one H200 (batch 1, bf16).
Training cost is for one DROID run.}
\label{tab:efficiency}
\small
\begin{tabular}{@{}lcccc@{}}
\toprule
Predictor & Params (M) & GFLOPs / 10 steps & Latency (ms) & Train GPU-h \\
\midrule
AR-TF / AR  & 21.70 & \pend & \pend & \pend \\
Direct      & 21.70 & \pend & \pend & \pend \\
\ours{}     & 21.75 & \pend & \pend & \pend \\
\bottomrule
\end{tabular}
\end{table}

\paragraph{Efficiency.} \Cref{tab:efficiency} compares parameters, compute per forecast step and
latency. Transport adds \pend{} parameters. Because \ours{} and Direct decode all $K$ steps in
parallel, a $K$-step forecast needs one encoder pass, whereas recursive predictors need $K$
(\pend{}).
